\documentclass[12pt]{article}
%friendlyname:Noɬdith;Owl
\usepackage{covington,graphicx,tipa,makeidx,fontspec}
\setmainfont[Mapping=tex-text]{Charis SIL}

\begin{document}

\title{NOŁDITH GIXIDHOY \newline (Owl Story)}
\author{Belle Deacon}
\date{}

\maketitle

Recorded by Karen MacPherson?  

Recording: anl0764b.wav. Length: 0:03:06; DX portion 0:02:29.

Told to John Deacon.

\bigskip 

BD: You know they said, long time ago, there was little
animals. They were, they turn into person, they said.

JD:  They were human one time.

BD: Human and then after that they turn into little like crow and
everything. And this little hawk,\footnote{owl. See next footnote.} she had four little young ones. Ah,
in Indian way you want?

\gll Noɬdith        tr'i'ne   yixidentl'iyo.
owl {it's said} {they're sitting in a line}
\glt `They say that some owls were sitting up high somewhere.'\footnote{In Osgood 1959:36,
\em niɬda\texttheta\ \em is identified as `owl-B', with
`longer and thicker feathers' than a `great horned owl'. In James Kari's
fieldwork with Cleo Fairbanks, \em niɬdith \em is
identified as `great gray owl'. However, the title of this story is translated
`hawk owl' in Deacon 1985, and that is how \em noɬdith \em is translated
here.}  
\glend

n

\gll Axaxiɬdik {xivogh dingit'a}           ɬogot. 
then      {she cares for them} truly
\glt `She (the mother) was pleased with her children .'
\glend

\gll Dits'inqay   niɬ'anh          ts'in', xits'in' ntthi'it'ox      ɬogot. 
{her children} {she's looking at} and    {to them}  {she's moving her head} truly
\glt `She nodded towards them.'
\glend

JD:  (laughs)

BD:  
\gll {``Aaa gila gila ginagheyo   ggiy ji"} xiɬne ɬogo  xiniɬ'anh hingo hinh.
{                  }   {she said to them} truly {she's looking at them} while she
\glt `...she sang to them as she looked at them.'
\glend

\gll Xits'i  nigidluq ts'i xiniɬ'anh             hingo.
{at them} {she's smiling} {}  {she's looking at them} while
\glt `She smiled at them as she looked at them.'
\glend

\gll Yiggi che giɬigginh cheche niɬ'anh.
that  again  {another one}       again  {she's looking}
\glt `She looked at another one of them again.' 
\glend

\gll {Aaa gila gila ginagheyo ggiy ji} xiɬdik
{} then
\glt `\ldots then'
\glend

\gll Uxuxuyiɬ ngu'odz  yit    tr'o'isr          ɬogot Yixgitsiy.
suddenly {from over} there {someone is coming} truly Crow
\glt `Suddenly Crow was coming over.'
\glend

\gll ``Isda,''          xivazrne.
{friend/cousin} {he said to them }
\glt ```Friend,'' he said to them.'
\glend

\gll ``Ngo  gits'i  {ximo'idendhen eden},'' xivazrne.
well wrongly {}  {he said to them}
\glt ```You don't speak to them the right way,'' he told (her).'\footnote{\em ximo eden \em apparently said in Holikachuk}
\glend

JD:  (laughs)

\gll ``Ithe'  gitth'in tisrgi ximdene ts'i'et,'' vazrne         ɬogo.
should leg      {big fat}    {you say to them}    {}     {he said to her} truly
\glt ```Why don't you call them big fat legs,'' he asked her.'\footnote{\em ximdene \em and \em ts'i'et \em in Holikachuk. DX words would be \em xividene \em and \em ts'i'at.}
\glend

That's up the river\footnote{in Holikachuk}, this crow. nnn ts'in' he don't no-- pay no attention.

\gll {A gila gila ginagheyo   go ɬe che} they sang. 
{            } 
\glt `\ldots, she sang.'
\glend

Nnn 

\gll {Pretty soon} ``ithe'  gimagg chux xividene  ts'i'at,'' vazrne      ɬogo.
{}                should eye    big  {tell them} {}    {he told her} truly
\glt `Pretty soon he told her, ``why don't you say `big eyes' to them,'' he asked her.' 
\glend

\gll Ngo uxuxiyiɬ yitongo   iɬt'et       ɬogo  gidile        {go dadz}  di'ne ts'in': {A gila gila ginagheyo   ggiy ji} \em {they said}
well  suddenly       meanwhile continuously truly {she's singing} {this way} {she's uttering}
\glt `All this time she kept on singing saying like this \ldots.'
\glend

\gll {`` `Ganhtse}    {ginan-ghidizriq '} xividene  anh,'' vazrne. 
{one's nose} {they're hooked/crooked} {tell them} he  {he told her}
\glt ```Tell them their noses are crooked,'' he said to her.' 
\glend

\gll Xiyiɬ xizronche yuxudz {yits'in' nughoɬtl'itthdi.}
then apparently        just   {she attacked him}
\glt `At that she attacked him.'
\glend

\gll Yuxudz, tthidixitadhit        yitots'in'.
just   {they started to fight} then
\glt `Then they started to fight.' 
\glend

\gll Eyggi yixgitsiy yuxudz {xitr'eɬtlux}     ɬogot.
that  raven     just   {she knocked down} truly 
\glt `She knocked that crow right down.' 
\glend

\gll ``Ggaq,'' didiyoq.
caw  {he uttered }
\glt ```Caw,'' he shouted.' 
\glend

\gll Che   diggandi'eɬtrit   cheche tthidanxitithidhit. 
again {he quickly got up again} again  {they started to fight again}
\glt `He quickly got back up and they started fighting again.' 
\glend
%BD says tthidanxitidhidhit

\gll Eygginh noɬdith cheche xitr'eɬtlux.
that    owl     again  {he knocked down}
\glt `And this time he knocked the owl down too.' 
\glend

\gll ``Ix,'' didiyoq. 
{} {she uttered}
\glt ```\em Ix,'' \em she shouted.'
\glend

\gll Che   tthidanxitithidhit    tux  ***   Yuxgitsiy che   xantr'etht'ix.
again {they start to fight again} when  raven     again {she knocked down}
\glt `They started fighting again and she flung Crow down again.' 
\glend
%BD says tthidanxitthit

\gll ``Ggaq,'' tr'i'ne ts'i yuxudz tr'ine--- tr'inliq'ak ngidiggi.
caw  {he says} and  just      {}    {he flew off}  above 
\glt `Saying ``caw'' he just flew away overhead.' 
\glend

\gll Ngo  yiggiy noɬdith ggoy  nginiggi yuxudz tthik'ugixenathtl'isr. 
well that   owl     young back     just   {they went into the woods}
\glt `Then the young owls all flew away into the woods.' 
\glend

\gll Ts'an ts'in' idixinlu'on'. 
then {}     {the end}
\glt `That's all.' 
\glend

This crow he fly away. That's the end.

\bigskip
Osgood, Cornelius. 1959. Ingalik mental culture. New Haven: Dept. of
Anthropology, Yale University. 

\end{document}